% Preamble shared by every model comparison document (docs/models/build.py pastes it into each generated .tex,
% so a document is one self-contained file). The same families, colours and code style as the engine manual.

\usepackage{fontspec}
\usepackage{amsmath,amssymb}
\usepackage[margin=1in,headheight=14pt]{geometry}
\usepackage{xcolor}
\usepackage{graphicx}
\usepackage{float}  % [H]: a figure stays where the text refers to it
\usepackage{microtype}
\usepackage[nobottomtitles*]{titlesec}
\renewcommand{\bottomtitlespace}{0.12\textheight}  % a heading this close to the foot of a page moves to the next one
\usepackage{needspace}
\usepackage{enumitem}
\usepackage{booktabs,array,xltabular}
\usepackage{listings}
\usepackage{fancyhdr}
\usepackage{parskip}
\usepackage{hyperref}

\definecolor{accent}{HTML}{9A5A00}
\definecolor{ink}{HTML}{1C1F22}
\definecolor{muted}{HTML}{5D636A}
\definecolor{codebg}{HTML}{F3F4F1}
\definecolor{passgreen}{HTML}{2F6B3F}

\hypersetup{colorlinks=true, linkcolor=ink, urlcolor=accent, citecolor=ink, pdfauthor={Anisotropy AI}}

\titleformat{\section}{\sffamily\Large\bfseries\color{ink}}{\color{accent}\thesection}{0.7em}{}
\titleformat{\subsection}{\sffamily\large\bfseries\color{ink}}{\color{accent}\thesubsection}{0.6em}{}
\titleformat{\subsubsection}{\sffamily\normalsize\bfseries\color{ink}}{\thesubsubsection}{0.5em}{}
\titlespacing*{\section}{0pt}{2.2ex plus 1ex}{1.2ex}
\titlespacing*{\subsection}{0pt}{2ex plus 1ex}{0.8ex}

\setlist{nosep, leftmargin=1.4em, topsep=0.3ex}
\setlength{\emergencystretch}{3em}
\newcolumntype{Y}{>{\raggedright\arraybackslash}X}
\renewcommand{\arraystretch}{1.25}

\lstdefinestyle{engine}{
  language=Python, basicstyle=\ttfamily\footnotesize, columns=fullflexible, keepspaces=true,
  numbers=left, numberstyle=\tiny\color{muted}, numbersep=7pt,
  backgroundcolor=\color{codebg}, frame=leftline, framerule=2pt, rulecolor=\color{accent},
  xleftmargin=2.8em, framexleftmargin=2.2em, aboveskip=0.4ex, belowskip=1.2ex,
  keywordstyle=\color{accent}\bfseries, commentstyle=\color{muted}\itshape,
  stringstyle=\color{passgreen}, showstringspaces=false, breaklines=true}
\lstnewenvironment{enginecode}[1][]{\lstset{style=engine,#1}}{}
% Plain text files (a spec, a record, a deck header): the same box without Python highlighting.
\lstdefinestyle{plain}{
  language={}, basicstyle=\ttfamily\footnotesize, columns=fullflexible, keepspaces=true,
  numbers=none, backgroundcolor=\color{codebg}, frame=leftline, framerule=2pt, rulecolor=\color{muted},
  xleftmargin=1.2em, framexleftmargin=0.8em, aboveskip=0.4ex, belowskip=1.2ex, breaklines=true}
\lstnewenvironment{plainfile}[1][]{\lstset{style=plain,#1}}{}

% \codefrom{function}{file}{lines}: caption above a code listing.
\newcommand{\codefrom}[3]{\par\needspace{6\baselineskip}\smallskip\noindent{\sffamily\footnotesize\color{muted}%
  Code: \textcolor{ink}{\textbf{#1}} in #2, lines #3}\par\nopagebreak}
% \filefrom{file}: caption above a listed file.
\newcommand{\filefrom}[1]{\par\smallskip\noindent{\sffamily\footnotesize\color{muted}File: \textcolor{ink}{#1}}\par\nopagebreak}
% \evidence{TAG}: the evidence tag after a value in the rebuild guide (INP, INP-derived, SPEC, CALC, CAE-CHECK, MANUAL p.N, GUI-SEEN, ASSUMED, OPEN).
\newcommand{\evidence}[1]{{\sffamily\scriptsize\color{accent}[#1]}}
% \verdict{text}: the boxed statement on the title page and at the head of the summary.
\newcommand{\verdict}[1]{\colorbox{codebg}{\parbox{\dimexpr\linewidth-2\fboxsep}{\small\color{ink}#1}}}
